\documentclass[aspectratio=169]{beamer}

\usepackage[utf8]{inputenc}
\usepackage[T1]{fontenc}
\usepackage{tikz}
\usepackage{booktabs}
\usepackage{amsmath}

\usetheme[noto, showmaxslides]{pureminimalistic}
\newcommand{\CourseCode}{MAS512}
\renewcommand{\pageword}{Page}
\renewcommand{\logoheader}{\vspace{1.5em}}

\title[\CourseCode]{Assignment 1 -- Q2, Q4, Q6}
\newcommand{\undertittel}{Object tracking, motor fault classification and autoencoders}
\author{Richard Rumpelt \and Javier Comes Tello \and Tobias Brand}
\institute{University of Agder}
\date{\today}

\graphicspath{{Figures/}}
\newcommand{\fig}[2][]{\includegraphics[#1]{#2}}
\newcommand{\good}[1]{\textbf{#1}}

\begin{document}

\makestartpage
\maketitle

\begin{frame}[plain, noframenumbering]{Outline}
    \tableofcontents
\end{frame}

% =====================================================================================
\section{Q2 -- Object Tracking}
\makechapterpage{Q2 -- Object Tracking}{5 points}

\begin{frame}{Objective \& Method}
    \begin{itemize}
        \item \good{Objective:} track the swinging load in \texttt{test\_swingingload.avi} and show its trajectory as it swings
        \item \good{Method:}
        \begin{itemize}
            \item Detector: the YOLO model trained in Q1 (\texttt{best.pt})
            \item \texttt{model.track(frame, persist=True)} on every frame -- keeps the same object identity across frames
            \item Bounding-box center $(c_x, c_y)$ extracted for every frame the load is detected in
            \item Plot all centers as a scatter of pixel coordinates $X$ vs.\ $Y$
        \end{itemize}
    \end{itemize}
\end{frame}

\begin{frame}{Result}
    \centering
    \fig[width=.62\linewidth]{q2_trajectory}
    \vfill
    {\small Load detected and tracked in \good{all 300} test-video frames -- the pendulum-like swing arc is clearly visible}
\end{frame}

% =====================================================================================
\section{Q4 -- Classification}
\makechapterpage{Q4 -- Classification}{10 points}

\begin{frame}{Objective \& Data}
    \begin{itemize}
        \item \good{Objective:} determine motor state of health (healthy vs.\ faulty) and compare an own CNN against the baseline architecture given in class
        \item \good{Data:} \texttt{fault-classification-class} -- vibration scalograms (laser vibrometer $\rightarrow$ wavelet transform), $60\times175\times3$ px
        \item 180 measurements per class, each split into 5 time-window crops $\Rightarrow$ 900 images/class
        \item Fixed 80/20 split by measurement number (1--144 train, 145--180 val) -- all crops of one measurement stay together, no data leakage
    \end{itemize}
    \centering
    \small
    \begin{tabular}{lccc}
        \toprule
        & Training & Validation & Test \\
        \midrule
        Images & 1440 & 360 & 200 \\
        \bottomrule
    \end{tabular}
\end{frame}

\begin{frame}{Model 1 (own) vs.\ Model 2 (given baseline)}
    \centering
    \fig[width=.85\linewidth]{model_1_architecture}
    \vfill
    \begin{itemize}
        \footnotesize
        \item \good{Model 1 (custom, 27.9k params):} 3$\times$[Conv2D+MaxPool] (16$\to$32$\to$64), GlobalAveragePooling, Dense(64)
        \item \good{Model 2 (given, 105k params):} 5$\times$[Conv2D+MaxPool+ReLU] (8$\to$16$\to$64$\to$64$\to$64), Flatten(320), FC(64)
    \end{itemize}
\end{frame}

\begin{frame}{Training}
    \begin{columns}[c]
        \begin{column}{.5\linewidth}
            \centering
            \fig[width=\linewidth]{q4_model1_loss}
        \end{column}
        \begin{column}{.5\linewidth}
            \centering
            \fig[width=\linewidth]{q4_model2_loss}
        \end{column}
    \end{columns}
    \vfill
    {\small Adam, categorical cross-entropy, \texttt{EarlyStopping} on validation loss (patience 3, best weights restored)}
\end{frame}

\begin{frame}{Evaluation}
    \begin{columns}[c]
        \begin{column}{.5\linewidth}
            \centering
            \fig[width=\linewidth]{q4_model1_confusion}
        \end{column}
        \begin{column}{.5\linewidth}
            \centering
            \fig[width=\linewidth]{q4_model2_confusion}
        \end{column}
    \end{columns}
\end{frame}

\begin{frame}{Comparison \& Conclusion}
    \centering
    \small
    \begin{tabular}{lccccc}
        \toprule
        Model & Params & Inference & Accuracy & Precision & F1 \\
        \midrule
        Model 1 (custom)   & \good{27{,}874}  & 0.76\,ms & 0.985 & 0.985 & 0.985 \\
        Model 2 (baseline) & 105{,}202 & 0.76\,ms & 0.985 & 0.985 & 0.985 \\
        \bottomrule
    \end{tabular}
    \vfill
    \begin{itemize}
        \item Both models are close to perfect -- the two classes are visually very distinct in the scalograms
        \item Model 1 reaches the \good{same accuracy with 3.8$\times$ fewer parameters}, mainly from GlobalAveragePooling instead of Flatten
    \end{itemize}
\end{frame}

% =====================================================================================
\section{Q6 -- Autoencoder}
\makechapterpage{Q6 -- Autoencoder}{10 points}

\begin{frame}{Objective \& Method}
    \begin{itemize}
        \item \good{Objective:} compress each scalogram to a latent vector and reconstruct it, for 10 random images per class
        \item \good{Method:} dense autoencoder (matches the architecture shown in class)
    \end{itemize}
    \centering
    \[
        \underbrace{60\times175\times3}_{31{,}500}
        \;\xrightarrow{\text{Dense}}\; 256
        \;\xrightarrow{\text{Dense}}\; \underbrace{32}_{\text{latent } z}
        \;\xrightarrow{\text{Dense}}\; 256
        \;\xrightarrow{\text{Dense}}\; \underbrace{60\times175\times3}_{31{,}500}
    \]
    \begin{itemize}
        \item Trained unsupervised: target = input itself, loss = MSE (reconstruction error)
    \end{itemize}
\end{frame}

\begin{frame}{Result -- Healthy (00\_15)}
    \centering
    \fig[width=.95\linewidth]{q6_healthy}
\end{frame}

\begin{frame}{Result -- Faulty (10\_15)}
    \centering
    \fig[width=.95\linewidth]{q6_faulty}
\end{frame}

\begin{frame}{Conclusion}
    \centering
    \begin{tabular}{lc}
        \toprule
        Class & Mean reconstruction error (MSE) \\
        \midrule
        Healthy (00\_15) & \good{0.0022} \\
        Faulty (10\_15)  & 0.0028 \\
        \bottomrule
    \end{tabular}
    \vfill
    \begin{itemize}
        \item Both classes reconstruct well -- the 32-value latent vector preserves the overall band pattern
        \item Faulty images have a slightly higher error: the periodic peak pattern carries more high-frequency detail than the smooth healthy band, making it harder to compress into 32 numbers
    \end{itemize}
\end{frame}

\end{document}
